% p162_1; source PDF page 162; manually verified editable table.
\begin{tikzpicture}[x=1cm,y=-1cm,inner sep=0pt,outer sep=0pt]
\path[use as bounding box] (0,0) rectangle (16.000,4.410);
\draw[line width=.4pt] (0.000,0.000) rectangle (2.200,1.205);
\node[align=center,text width=1.940cm,font={\linespread{1}\fontsize{10}{13.50}\selectfont\bfseries}] at (1.100,0.603) {作为义务人（甲）};
\draw[line width=.4pt] (2.200,0.000) rectangle (4.600,1.205);
\node[align=center,text width=2.140cm,font={\linespread{1}\fontsize{10}{13.50}\selectfont\bfseries}] at (3.400,0.603) {非作为义务人（乙）};
\draw[line width=.4pt] (4.600,0.000) rectangle (10.200,1.205);
\node[align=center,text width=5.340cm,font={\linespread{1}\fontsize{10}{13.50}\selectfont\bfseries}] at (7.400,0.603) {范例};
\draw[line width=.4pt] (10.200,0.000) rectangle (16.000,1.205);
\node[align=center,text width=5.540cm,font={\linespread{1}\fontsize{10}{13.50}\selectfont\bfseries}] at (13.100,0.603) {结论};
\draw[line width=.4pt] (0.000,1.205) rectangle (2.200,4.410);
\node[align=left,text width=1.940cm,font={\linespread{1}\fontsize{10}{13.50}\selectfont}] at (1.100,2.808) {实行};
\draw[line width=.4pt] (2.200,1.205) rectangle (4.600,4.410);
\node[align=left,text width=2.140cm,font={\linespread{1}\fontsize{10}{13.50}\selectfont}] at (3.400,2.808) {帮助、教唆};
\draw[line width=.4pt] (4.600,1.205) rectangle (10.200,4.410);
\node[align=left,text width=5.340cm,font={\linespread{1}\fontsize{10}{13.50}\selectfont}] at (7.400,2.808) {1. 甲不想履行扶养义务，要将老父亲丙遗弃，让乙帮忙。乙照办。2. 甲不慎导致丙受伤，想救丙。乙教唆甲不要救。甲便不救};
\draw[line width=.4pt] (10.200,1.205) rectangle (16.000,4.410);
\node[align=left,text width=5.540cm,font={\linespread{1}\fontsize{10}{13.50}\selectfont}] at (13.100,2.808) {甲构成不作为犯的实行犯，乙构成帮助犯、教唆犯（作为义务人是一种身份，乙没这个身份，故不可能构成不作为犯的实行犯）};
\end{tikzpicture}
